\section{Revised next-stage routes}
\label{sec:revised-routes}

\subsection{Route A: anchored primitive-plane variation}

For one row pair and one full same-residual anchor \(\mathfrak a\), let
\(U_{\mathfrak a},V_{\mathfrak a}\) be the actual primitive-coordinate
extents of the aggregated literal kernel
\(\cK_{\mathfrak a}^{\circ}\) from
\eqref{eq:tpc72-aggregate-erasure-kernel}.  Define the exact Abel cost
\begin{equation}
 \mathfrak A_{\mathfrak a}(D)
 :=
 \sum_{d\le D}
 \mathfrak M_d(U_{\mathfrak a}/d)
 \mathfrak M_d(V_{\mathfrak a}/d)
 \cV_\square((\cK_{\mathfrak a}^{\circ})^{[d]})
 +
 \cT_D(\cK_{\mathfrak a}^{\circ}).
 \label{eq:tpc72-Abel-cost}
\end{equation}
Define the anchored plane mass, with absolute values taken after the exact
\(q\)- and ray-aggregation, by
\begin{equation}
 \mathfrak W_{\mathfrak a}
 :=
 \sum_{u,v\ge1}
 |\cK_{\mathfrak a}^{\circ}(u,v)|.
 \label{eq:tpc72-anchored-plane-mass}
\end{equation}
Because \(\cK_{\mathfrak a}^{\circ}\) already contains the normalized
four-corner amplitudes,
\[
 |\mathfrak X_{\mathfrak a}|
 \le
 \mathfrak W_{\mathfrak a}.
\]
This mass can be smaller than the sum of absolute pre-aggregation ray
atoms, and it is not automatically the global trace or raw atomic mass.
A first useful arithmetic target is
\begin{equation}
 \mathfrak A_{\mathfrak a}(D_X)
 \le
 X^{-\tau_{\rm plane}+o(1)}\mathfrak W_{\mathfrak a}
 \label{eq:tpc72-plane-saving}
\end{equation}
uniformly over every active anchor, pair, block, and key in a declared
scope.  The physical conclusion comes only after the complete anchor and
row sum in \eqref{eq:tpc72-correct-trace-majorant}.
If \(D_X\) grows, \eqref{eq:tpc72-plane-saving} additionally requires a
proved restricted-Mertens estimate uniform throughout \(d\le D_X\).
With only fixed-\(d\) input, the safe formulation fixes \(D\), takes the
physical scale limit, and only then lets \(D\to\infty\).

This target must be attacked in the following order:
\begin{enumerate}[label=\textup{(A\arabic*)},leftmargin=3.6em]
 \item compute the actual support and form the literal \(q\)-aggregate
       \(\cK_{\mathfrak a}^{\circ}\) before absolute values;
 \item compute its zero-extended mixed variation, including bulk, edge, and
       corner jumps;
 \item estimate the exact tail
       \(\cT_D(\cK_{\mathfrak a}^{\circ})\); use
       \eqref{eq:tpc72-ambient-tail} only when that specialization is sharp
       enough;
 \item choose a fixed cutoff or justify a growing \(D_X\) against the
       physical mass rather than ambient area;
 \item test whether the resulting saving survives complete trace
       aggregation.
\end{enumerate}
If a sharp lower obstruction proves that the exact mass-normalized cost in
\eqref{eq:tpc72-Abel-cost} is necessarily trivial on a nonnegligible anchor
family, this Abel certificate is closed for that carrier.  Failure of the
ambient upper bound alone closes only the ambient-tail specialization, and
failure to prove the sufficient inequality is an unresolved gate rather
than a negative theorem.

\subsection{Route B: build a genuine operator carrier}

Equation \eqref{eq:tpc72-plane-saving} does not give an entrywise pair
bound.  To invoke the TPC--68 pair/star machinery, one must first construct
a direct primitive or other arithmetic representation of the two-corner
entry \eqref{eq:tpc72-two-corner-entry}.  If that new representation gives
an entrywise saving \(X^{-\tau_{\rm ent}+o(1)}\) and the actual pair graph
has degree exponent \(\delta\), the TPC--68 pair route requires
\begin{equation}
 \tau_{\rm ent}\ge\delta
 \quad\text{for a subpower norm,}
 \qquad
 \tau_{\rm ent}>\delta
 \quad\text{for decay.}
 \label{eq:tpc72-pair-threshold}
\end{equation}
If square summation improves the local estimate to star exponent \(\sigma\),
the threshold becomes
\begin{equation}
 \sigma\ge\delta/2
 \quad\text{or}\quad
 \sigma>\delta/2.
 \label{eq:tpc72-star-threshold}
\end{equation}
Below these thresholds the corresponding certificate supplies neither a
subpower bound nor decay by itself.  It does not kill the whole operator
route: a positive norm exponent may still be absorbed by the mass-degree
gate \eqref{eq:tpc72-zero-cost-reassembly}.

A second bridge may instead prove row-local anchored bounds that dominate
the actual stars without first identifying every entry.  This requires a
new row-local crosswalk and cannot be inferred from the global trace.
A third bridge is to estimate complete decorated moments with constants
uniform in the growing length.  The moment option requires a new crosswalk
for every decorated sector actually used; the first mixed trace cannot
stand in for the full moment.

The global trace \(\tr(ER)\) is useful as a diagnostic of the linear sector,
but it is not an operator aggregation target.  The next paper must therefore
choose explicitly among a direct two-corner entry route, a row-local
anchored-star bridge, and a full-moment route.  A fixed-row star or
reducible-safe Perron majorant becomes available only after a literal
carrier dominating that object has been constructed.

\subsection{Route C: represented-frame closure in parallel}

The primitive-plane branch should not monopolize the program.  Independently,
one needs:
\begin{enumerate}[label=\textup{(R\arabic*)},leftmargin=3.6em]
 \item a terminal-candidate lower theorem giving
       \eqref{eq:tpc72-uniform-activation-target} or its distinguished
       analogue;
 \item a physical represented Gram lower bound, through determinant mass,
       a comparison matrix, or a different exact certificate;
 \item an explicit physical label-erasure or fusion-angle estimate.
\end{enumerate}
The determinant route is alive precisely because its conversion is exact,
but it needs actual matching, pivot, and cycle-activity input.  If the worst
actual component determinant mass is proved superpolynomially small, that
closes this certificate at endpoint scale while leaving other represented
certificates open.

\subsection{Route D: final mass-degree and shorting merge}

Only after Routes A--C give compatible inputs should one test
\begin{equation}
 \sigma_M>\nu_M+\lambda_{\rm rep}
 \label{eq:tpc72-final-mass-degree}
\end{equation}
and
\begin{equation}
 \|C_N\|=o(1).
 \label{eq:tpc72-final-nuisance}
\end{equation}
Both must use the same reference energy, carrier, support, coefficient
scope, and fixed shift.  The endpoint ledger is then evaluated once, with
boundary and tail attached last.

\subsection{Fallbacks if the anchored trace door closes}

If the physical primitive-plane variation is too large, or if its residual
boundary is frozen on the actual support, the remaining signed alternatives
are:
\begin{enumerate}[label=\textup{(\roman*)},leftmargin=3em]
 \item higher one-erasure polygons;
 \item cancellation from literal complex phases;
 \item a direct shorted reassembly estimate; or
 \item a coefficient-specific distinguished theorem with its scope stated
       explicitly.
\end{enumerate}
The complete quadratic sector is not a fallback source of negative
compensation by \eqref{eq:tpc72-quadratic-positive}.

\subsection{Stage schedule}

A disciplined next block can be organized as follows.
\begin{enumerate}[label=\textbf{Stage \arabic*.},leftmargin=4.5em]
 \item Anchored physical variation census and boundary decomposition.
 \item A uniform or obstructed first-mixed-trace estimate.
 \item A direct two-corner, row-local anchored-star, or complete-moment
       operator crosswalk.
 \item Pair-to-star/Perron aggregation or uniform moment control, according
       to the carrier established in Stage 3.
 \item Parallel terminal activation and represented determinant work.
 \item Missing-mass and nuisance-shorting comparison.
 \item Boundary, tail, and strict endpoint audit.
\end{enumerate}
At every stage, success means either a literal exponent improvement, a
necessary interface closure, or an exact actual-carrier obstruction.  A
model refinement without one of these outputs should not be reported as
movement toward the parity endpoint.
